\section{Related Work}
This section organizes prior work on mobility-on-demand systems by the Markov decision process (MDP) components used in the literature: \textit{State}, \textit{Action}, \textit{Transition}, and \textit{Reward}. For each component we compare representative approaches, highlighting similarities and differences in modeling choices, level of abstraction, and use of learning methods.

\subsection{State}
Works differ substantially in how they represent the system state, from macroscopic network-flow summaries to fine-grained mesoscopic or agent-level descriptions. Network- or zone-level aggregate states (idle vehicles, waiting passengers, OD rates) are common in optimization and MPC approaches such as the spatiotemporal pricing and dynamic-programming framework \citep{spatiotemporal_pricing} and the spatial-pricing equilibrium models \citep{spatial_pricing,geometric_matching}. These macroscopic states enable tractable dynamic programming or equilibrium analysis but typically assume exogenous travel times and ignore endogenous congestion effects (or model them via simple speed functions) \citep{spatial_pricing,spatiotemporal_pricing}.

Mesoscopic and agent-based simulators represent more detailed state that includes vehicle-level attributes, region-level counts, and traveler requests. Several ABMS and simulation-driven studies encode per-vehicle and per-request state (vehicle location, status, queued requests) and use these richer states for empirical evaluation: hybrid HV/AV simulations \citep{hybrid_operations}, multi-operator ABMs \citep{modeling_the_competition}, and broker/ridepooling games \citep{competition_and_cooperation}. Papers that apply learning methods tend to adopt intermediate representations that balance expressiveness with scalability: graph-based states that summarize node-level supply/demand and spatial relations are used in GNN/graph-RL works \citep{graph_meta,learning_based}, while mesoscopic hex-grid or zone vectors appear in MADRL and two-sided DRL studies \citep{competitive_pricing,two_sided}.

In short, optimization-oriented work prefers aggregated states for tractability \citep{spatiotemporal_pricing,geometric_matching,spatial_pricing}, simulation and ABM literature use richer agent-level states for realism \citep{hybrid_operations,modeling_the_competition,competition_and_cooperation}, and learning papers compromise with structured representations (graphs, zonal vectors) for scalability and transfer \citep{graph_meta,learning_based,competitive_pricing,two_sided}.

\subsection{Action}
Action spaces reflect the control authority considered: platform-level continuous controls (prices, commission rates, rebalancing flows), vehicle-level discrete relocation/dispatch decisions, and agent-level acceptance or offline choices. Centralized platform controls (pricing and rebalancing) are the focus of network-flow and pricing optimization papers \citep{spatiotemporal_pricing,geometric_matching,spatial_pricing}, which treat actions as continuous pricing or flow variables amenable to dynamic programming or convex/dual decomposition techniques.

Agent- or vehicle-level discrete actions (accept/decline, relocate, pick up, offer construction) arise in ABM and game-based studies \citep{competition_and_cooperation,modeling_the_competition,hybrid_operations}. Decentralized RL approaches likewise model per-vehicle actions (rebalance, pick up, do nothing) or driver decisions, as in SAMoD and two-sided DRL where vehicles/drivers learn local policies \citep{samod,two_sided}. Multi-agent and graph-based RL formulations often combine higher-level continuous actions (desired distribution shares, price coefficients) with lower-level solvers that convert these signals to executable flows or prices \citep{learning_based,graph_meta,competitive_pricing}.

Thus the literature spans from continuous, centralized controls for optimization (pricing/flow) to discrete, decentralized controls for operational decisions, with hybrid formulations common in MARL and GNN-based control where high-level continuous actions are post-processed by specialized solvers \citep{learning_based,graph_meta,competitive_pricing}.

\subsection{Transition}
Transition models determine how actions and stochastic demand produce next states. Papers focused on analytical insight often propose stylized or relaxed transition dynamics: geometric matching and matching-time approximations in \citep{geometric_matching} or static equilibrium constraints in \citep{spatial_pricing} enable analytic characterization and welfare comparisons. The macroscopic network-flow work explicitly derives differential equations for aggregate dynamics and uses relaxations to enable decomposition \citep{spatiotemporal_pricing}.

By contrast, many empirical and learning studies rely on bespoke simulators to capture operational detail: mesoscopic simulators with routing and matching (KM) are central to competitive pricing, hybrid operations, and multi-operator ABMs \citep{competitive_pricing,hybrid_operations,modeling_the_competition}. These simulators often integrate routing engines (OSRM or Dijkstra variants), matching heuristics or the Hungarian algorithm, and vehicle state machines. Game-based ridepooling and broker studies also implement insertion heuristics and ILP re-optimizations for realistic ridepool schedules \citep{competition_and_cooperation}. Learning methods typically train policies in such simulators (simulator-in-the-loop) or use the simulator as the environment for RL; meta-RL and transfer works assume simulated task distributions across cities and rely on simulators for transitions \citep{graph_meta,learning_based}.

In short, analytic-transition papers favor tractable, stylized models for insight and bounds \citep{geometric_matching,spatiotemporal_pricing,spatial_pricing}, while empirical and RL works rely on richer simulator-driven transitions to capture matching, routing and operational stochasticity \citep{competitive_pricing,hybrid_operations,modeling_the_competition,competition_and_cooperation,graph_meta}.

\subsection{Reward}
Objectives range from explicit platform profit maximization and welfare metrics to heuristic or sparse rewards used in decentralized RL. Many optimization and economic-model papers formulate platform profit or social welfare objectives directly (revenue minus costs, consumer/producer surplus) and analyze optimal pricing or regulatory interventions \citep{spatiotemporal_pricing,geometric_matching,spatial_pricing}. ABM/game studies evaluate operator profit, pooling efficiency and user-centered KPIs but do not always optimize via RL; rather they study strategic interactions or broker policies through experiments \citep{competition_and_cooperation,modeling_the_competition}.

Reinforcement-learning studies cast objectives as cumulative profit or surrogate performance signals. MADRL and multi-operator RL works typically use operator profit as the reward (revenue minus operational costs) \citep{competitive_pricing,learning_based,graph_meta}, while decentralized vehicle-level RL sometimes employs sparse or service-focused rewards (e.g., reward on successful passenger serving) to encourage occupancy and service \citep{samod}. Two-sided DRL frameworks include multi-headed objectives that combine platform profit, commission collection and service fulfillment metrics to balance conflicting goals between a supervisor and many drivers \citep{two_sided}.

Overall, optimization literature uses explicit economic objectives and welfare metrics for prescriptive analysis \citep{spatiotemporal_pricing,geometric_matching,spatial_pricing}, ABM/game work reports KPI-driven evaluation \citep{competition_and_cooperation,modeling_the_competition}, and RL work typically optimizes cumulative profit or service-related rewards, with some methods using sparse or shaped reward signals for decentralized agents \citep{competitive_pricing,learning_based,samod,two_sided}.

\subsection{Takeaways}
The reviewed works form a spectrum from analytic, optimization-driven models with aggregated states and closed-form or relaxed transitions to simulator-based, data-driven ABMs and RL studies with detailed vehicle-level states and complex transition dynamics. Graph-structured representations and hybrid action abstractions are emerging as effective compromises for scaling learning to urban-scale problems while preserving spatial structure and transferability \citep{graph_meta,learning_based,competitive_pricing}. For practitioners, the trade-off is clear: analytic models provide insight and bounds; simulators and learning enable operational realism and policy discovery but require careful calibration and computational resources.

% End of related work
